\section{Evaluación del modelo}

\subsection{La importancia de la evaluación}

El expositor admite que la evaluación de los LLM es uno de los mayores problemas de todo el campo: «en realidad no sabemos muy bien qué estamos haciendo». Pero la evaluación sigue siendo fundamental, porque la esencia del Post-Training es la \textbf{optimización}, y si la métrica de evaluación no es la correcta, la dirección de la optimización se desvía.

\subsection{Método de evaluación uno: Automated Benchmarks}

\textbf{Principio}: dado un conjunto de preguntas (como las de opción múltiple de MMLU), el modelo elige una respuesta y se le da una calificación con métricas como la exactitud.

\begin{table}[H]
\centering
\caption{Benchmarks automatizados comunes}
\begin{tabular}{lll}
\toprule
\textbf{Benchmark} & \textbf{Dominio} & \textbf{Métrica} \\
\midrule
MMLU & Conocimiento general (57 disciplinas) & Exactitud \\
GSM8K & Razonamiento matemático de primaria & Exactitud \\
MATH & Matemáticas de competencia & Exactitud \\
HumanEval & Generación de código & pass@k \\
IFEval & Instruction Following & Tasa de cumplimiento de restricciones \\
\bottomrule
\end{tabular}
\end{table}

\textbf{Ventajas}: escalable, de bajo costo, reproducible, se puede diseñar para tareas específicas.

\textbf{Desventajas}: no refleja el uso real del modelo; el usuario \textbf{conversa} con el modelo, no responde preguntas de opción múltiple.

\subsection{Método de evaluación dos: Human Evaluation (evaluación humana)}

\textbf{Principio}: evaluadores humanos eligen la mejor respuesta entre las de dos modelos anónimos (como en Chatbot Arena).

\textbf{Ventajas}:
\begin{itemize}
    \item Refleja directamente la preferencia humana, que es el objetivo último de optimización del Post-Training
    \item Permite controlar con precisión la dimensión evaluada (por ejemplo, centrarse solo en la toxicidad o solo en la exactitud)
    \item Riesgo bajo de contaminación de datos
\end{itemize}

\textbf{Desventajas}:
\begin{itemize}
    \item Es sumamente difícil de escalar, con costo alto y mucho tiempo
    \item Los evaluadores humanos \textbf{tienen sesgos importantes}: prefieren respuestas más largas y un tono más seguro, aunque el contenido sea incorrecto
\end{itemize}

\begin{warningbox}{La preferencia humana no se correlaciona con los benchmarks automatizados}
El expositor mostró un hallazgo clave: la clasificación de Chatbot Arena tiene una correlación muy baja con benchmarks automatizados como MMLU, GSM8K y HumanEval. Un modelo con puntuación alta en MMLU puede tener un desempeño muy pobre en la preferencia humana, y viceversa. Por eso, \textbf{es necesario usar ambos tipos de evaluación} a la vez para entender de manera integral la capacidad del modelo.
\end{warningbox}

\subsection{Método de evaluación tres: LLM-as-Judge}

Para encontrar un punto medio entre la escalabilidad y la preferencia humana, se puede usar un LLM en lugar de una persona para evaluar:

\begin{itemize}
    \item \textbf{Un solo Judge}: un LLM potente (como GPT-4) califica las respuestas
    \item \textbf{Varios Judge (Jury)}: un jurado formado por varios LLM vota para decidir cuál es mejor, lo que resulta más robusto y confiable
\end{itemize}

\textbf{Ventajas}: se correlaciona mucho con la preferencia humana, puede manejar tareas complejas de formato libre y da retroalimentación directa.

\textbf{Desventajas}: el LLM judge también tiene sesgos propios (y son muy parecidos a los sesgos humanos), por lo que aún se necesita una pequeña cantidad de evaluación humana para verificar la consistencia del judge.

\subsection{Recomendaciones para construir un sistema de evaluación propio}

\begin{enumerate}
    \item \textbf{Empezar lo antes posible}: diseñar la evaluación antes del fine-tuning, de forma similar al desarrollo guiado por pruebas
    \item \textbf{Iterar de manera continua}: el conjunto de datos de evaluación no es algo fijo, hay que corregirlo constantemente según el desempeño del modelo
    \item \textbf{Usarlos de manera combinada}: benchmark automatizado + LLM-as-Judge + una pequeña cantidad de evaluación humana
    \item \textbf{Comparar de forma transversal}: no solo comparar las distintas versiones del propio fine-tune, sino también contrastar con otros modelos y arquitecturas
\end{enumerate}

\begin{knowledgebox}{Desarrollo iterativo guiado por la evaluación}
El expositor enfatiza que la distribución del tiempo en el Post-Training debería ser, aproximadamente: un tercio en la preparación de datos, un tercio en el entrenamiento y un tercio en la evaluación. La evaluación no es una verificación posterior, sino el elemento central que atraviesa todo el proceso. El resultado de cada ronda de evaluación guía directamente la mejora de los datos y el entrenamiento de la siguiente ronda.
\end{knowledgebox}

\subsection{Resumen de la sección}

Existen principalmente tres métodos de evaluación de LLM: el benchmark automatizado (escalable, pero no refleja el uso real), la evaluación humana (la más cercana al objetivo, pero difícil de escalar) y LLM-as-Judge (una solución intermedia). Los tres se complementan y deben usarse de manera combinada. La baja correlación entre la preferencia humana y los benchmarks automatizados es un hecho clave que hay que reconocer al diseñar un sistema de evaluación.

\newpage
